\section{Experimental Setup and Results}
\label{results}

 Our experimental set-up is as follows:
\begin{enumerate}
    \item We start with a benchmark of coding problems 
    defined by both natural language and formal Verus specifications. 
    We consider these formal Verus specifications to be the ``gold-standard" 
    which ultimately defines the coding problem.
    \item We provide a model with just the natural language specifications, 
    and ask it to generate 
    (A) unit tests to check if an implementation satisfies the specification and 
    (B) a formal Verus specification. 
    Note that this generated spec may differ from the original gold-standard spec.
    \item For both (A) and (B), we evaluate the two models by checking if the code produced also satisfies
    the gold-standard spec. 
    If it does not, we interpret this to mean that the model successfully ``hacked''
    the reward mechanism. 
    We hypothesize that we will see more cases 
    when the produced code passes the unit tests but does not satisfy the gold-standard spec,
    compared to using generated specs as the reward mechanism.
\end{enumerate}
\begin{figure}[h]
\centering
\tikzset{
  box/.style={rectangle, rounded corners=3pt, draw, align=center,
              minimum width=2.0cm, minimum height=0.65cm, font=\small},
  goldbox/.style={box, fill=yellow!50, draw=orange!90!black, very thick},
  genbox/.style={box, fill=blue!30, draw=blue!80!black, very thick},
  rlbox/.style={box, fill=green!20, draw=green!60!black, very thick},
  evalbox/.style={box, fill=red!30, draw=red!80!black, very thick},
  arr/.style={-{Stealth[length=4pt]}, thick},
}
\begin{tikzpicture}[node distance=0.3cm and 1.0cm]

  % Top: benchmark input
  \node[goldbox, minimum width=3.4cm] (bench)
    {Benchmark\\{\footnotesize NL spec + gold Verus spec}};

  % NL spec fed to generator
  \node[rlbox, below=0.45cm of bench] (gen)
    {LLM Generator};
  \draw[arr] (bench) -- node[right, font=\footnotesize]{NL spec} (gen);

  % Two branches
  \node[genbox, below left=-.3cm and 1.2cm of gen]  (tests)
    {(A) Unit Tests};
  \node[genbox, below right=-.3cm and 1.2cm of gen] (spec)
    {(B) Generated Verus Spec};
  \draw[arr] (gen) -- (tests);
  \draw[arr] (gen) -- (spec);

  % RL training nodes
  \node[rlbox, below=0.35cm of tests] (rlA)
    {LLM Implementor\\{\footnotesize goal: pass unit tests}};
  \node[rlbox, below=0.35cm of spec]  (rlB)
    {LLM Implementor\\{\footnotesize goal: satisfy Verus spec}};
  \draw[arr] (tests) -- (rlA);
  \draw[arr] (spec)  -- (rlB);

  % Code outputs
  \node[box, below=0.35cm of rlA] (codeA) {Code A};
  \node[box, below=0.35cm of rlB] (codeB) {Code B};
  \draw[arr] (rlA) -- (codeA);
  \draw[arr] (rlB) -- (codeB);

  % Evaluation node
  \node[evalbox, minimum width=3.4cm,
        below=0.6cm of $(codeA)!0.5!(codeB)$] (eval)
    {Evaluate \\{\footnotesize reward hacking $\Leftrightarrow$ gold Verus spec not satisfied}};
  \draw[arr] (codeA) -- (eval);
  \draw[arr] (codeB) -- (eval);

  % Gold spec dashed arrow going around the right side
  \draw[arr, dashed, draw=orange!80!red!90!black, very thick] (bench.east)
        -- ++(4.2,0)
        |- (eval.east)
        node[pos=0.25, right, font=\small\bfseries, text=orange!80!red!90!black]{gold spec};

\end{tikzpicture}
\caption{Experimental setup. Both branches receive only the natural language specification as input. The dashed arrow indicates the gold-standard Verus specification used solely for final evaluation.}
\label{fig:setup}

\end{figure}

% Note: it's possible that the code produced by the models does satisfy the gold-standard specification,
%     but Verus rejects it due to insufficient proof annotations. 
%     If we find this to be the case, 
%     we would add a step between (3) and (4) which asks a model to add only proof annotations to the produced code,
%     to help Verus see that the implementation does satisfy the specification.

\textbf{Benchmarks}: We will evaluate this on the human eval Verus
benchmark~\cite{alphaverus}, which consists of around 100 coding problems with
both natural language specifications, formal Verus specifications, and a Rust
implementation. This is the best starting point: it is the largest known Verus
benchmark set with both natural language and formal specifications. Time
permitted, we could explore other benchmark sets.